\clearpage
\fancyhead[L]{\small TABLES DE LA LUNE.}
\SourceLeaf{645}
\begin{center}\small Göttingische gelehrte Anzeigen. Stück 23. Seite 225..227. 1825 Februar 7.\\[3pt]\rule{.72\linewidth}{.3pt}\end{center}

\emph{Tables de la lune, formées par la seule théorie de l'attraction, et suivant la division de la circonférence en \textup{400} degrés. Par le Baron \textup{\textsc{de Damoiseau}}, Lieutenant Colonel etc.} Paris 1824. Bei Bachelier 90 Seiten in 4.

So gross die Vollkommenheit ist, welche die Mondstafeln durch die Arbeiten von \textsc{Tobias Mayer}, \textsc{Bürg} und \textsc{Burckhardt} erhalten haben, so blieb ihr empirischer Ursprung doch gewissermaassen ein Vorwurf für die theoretische Astronomie. Die Genauigkeit, mit welcher man die Ungleichheiten der Mondsbewegung aus dem Princip der allgemeinen Schwere bisher ableiten konnte, blieb noch bedeutend hinter derjenigen zurück, welche man auf empirischem Wege erreicht hatte. Es war hier durch geschicktes Angreifen sehr verwickelter analytischer Rechnungen und durch sorgfältiges unermüdetes Durchführen derselben noch ein ehrenvoller Kranz zu erringen. Die Pariser Akademie munterte dazu auf, indem sie diese Rechnungen zum Gegenstand einer Preisaufgabe machte, und veranlasste dadurch die erfolgreichen Arbeiten von \textsc{Damoiseau}, \textsc{Carlini} und \textsc{Plana}. Hoffentlich werden wir bald die theoretischen Untersuchungen dieser verdienten Astronomen erhalten. Zuerst haben wir hier vor uns die auf \textsc{Damoiseau}'s Entwickelungen gegründeten Tafeln, die schon deren glücklichen Erfolg beurkunden.

Die Einrichtung dieser Tafeln weicht von der bisher gebrauchten darin ab, dass für die Länge des Mondes in der Ekliptik und für die Horizontal-Parallaxe sämmtliche Argumente der Zeit proportional sich ändern, während für die Breite die durch die Ungleichheiten schon verbesserte Länge zum Grunde liegt. Durch jene Einrichtung wird zwar die Convergenz bedeutend vermindert, und eine grössere Anzahl von Gleichungen erforderlich (sie sind für die Länge in 46 Tafeln gebracht, mit Einfluss zweier für die Störungen durch die Venus und den Jupiter), allein durch den ganz gleichförmigen Gang, welchen jetzt alle Rechnungen haben, wird dieses wohl ziemlich compensirt, obwohl das Zusammenfassen vieler kleiner Gleichungen in Tafeln mit doppelten Eingängen beim Gebrauch doch immer etwas unbequemes hat. Hätten die Tafeln mit doppelten Eingängen ganz vermieden werden sollen, so würde bei der Länge die Anzahl allen Tafeln auf 67 gestiegen sein. Die Benutzung der von \textsc{Carlini} in seinen Sonnentafeln gewählten Einrichtungen, die tägliche Aenderung der gleichförmig wachsenden Argumente zu ihrer Einheit anzunehmen, würde wohl den Gebrauch der Tafeln um vieles erleichtern.

Um die Genauigkeit, welche durch \textsc{Damoiseau}'s Mondstheorie erreicht ist, beurtheilen und mit der der frühern Tafeln vergleichen zu können, ist am Schluss des Werks die Vergleichung von 50 Beobachtungen, die während eines anderthalbjährigen Zeitraums 1802 und 1803 auf der Greenwicher Sternwarte angestellt sind, mit den \textsc{Bürg}'schen, \textsc{Burckhardt}'schen und \textsc{Damoiseau}'schen Tafeln beigefügt. Der blosse Anblick zeigt schon, dass die letzten die Beobachtungen nicht schlechter darstellen, als die erstern. Um indessen ein bestimmteres Urtheil fällen zu können, hat Ref. sich die Mühe gegeben, daraus den sogenannten mittlern Fehler abzuleiten, welcher sich
\begin{center}
\begin{tabular}{l|c|c}
bei den Tafeln von & in der Länge & in der Breite \\
\hline
Bürg & $5''80$ & $9''18$ \\
Burckhardt & $5.37$ & $7.82$ \\
Damoiseau & $5.23$ & $7.12$ \\
\end{tabular}
\end{center}
\noindent ergeben hat. So viel sich also aus dieser verhältnissmässig noch kleinen Anzahl von
Vergleichungen schliessen lässt, hat jetzt wirklich die Theorie dem Empirismus den Rang
abgewonnen. Ein nicht ganz \SourceLeaf{646} unbeträchtlicher Theil der Abweichungen mag übrigens wohl den Beobachtungen selbst zuzuschreiben sein; allein nachdem die Bahn einmal so weit gebrochen ist, scheint die Hoffnung nicht eitel zu sein, dass durch unermüdetes Fortschreiten auf demselben Wege die Vollkommenheit der Tafeln noch immer mehr vergrössert und die Darstellung der Mondsbewegung fast mit derselben Schärfe, deren die verfeinerte Beobachtungskunst fähig ist, erreicht werden wird.
